% ============================================================
% INTRODUCTION GÉNÉRALE
% ============================================================

\chapter*{Introduction Générale}
\addcontentsline{toc}{chapter}{Introduction Générale}
\markboth{Introduction Générale}{Introduction Générale}

\vspace{0.5cm}

L'essor fulgurant de l'intelligence artificielle générative constitue sans doute la mutation technologique la plus significative de cette décennie. Au c\oe ur de cette révolution se trouvent les modèles de langage de grande taille --- communément désignés par l'acronyme LLM (\textit{Large Language Models}) --- dont les capacités de compréhension et de génération de texte ont atteint un niveau de sophistication sans précédent. Les travaux fondateurs de Vaswani \textit{et al.}~\cite{vaswani2017} sur l'architecture Transformer, puis les recherches de Brown \textit{et al.}~\cite{brown2020} sur les propriétés émergentes des modèles à grande échelle, ont ouvert la voie à une nouvelle génération d'outils cognitifs : GPT-4 d'OpenAI~\cite{openai2023}, Claude d'Anthropic~\cite{anthropic2024}, Gemini de Google ou encore LLaMA de Meta~\cite{touvron2023}.

L'adoption de ces technologies par le monde industriel est désormais massive. Selon le rapport annuel de McKinsey \& Company publié en mars 2024, plus de 72\,\% des organisations à l'échelle mondiale ont intégré au moins une forme d'intelligence artificielle dans leurs processus opérationnels~\cite{mckinsey2024}. Bommasani \textit{et al.}~\cite{bommasani2022} qualifient les LLM de \og modèles fondamentaux \fg{} (\textit{foundation models}), soulignant leur caractère transversal et leur capacité à transformer simultanément de multiples secteurs d'activité.

\vspace{0.3cm}

Toutefois, cette adoption croissante engendre des problématiques majeures en termes de \textbf{gouvernance}, de \textbf{maîtrise des coûts} et de \textbf{traçabilité}. La consommation de tokens --- unité élémentaire de facturation des services LLM --- représente désormais un poste budgétaire significatif pour les entreprises, souvent opaque et difficilement prévisible. L'absence d'outils de supervision centralisés expose les organisations à des dépassements budgétaires non détectés, à des utilisations non conformes et à une incapacité d'identifier les anomalies de consommation en temps réel. Chen \textit{et al.}~\cite{chen2024llmops} identifient ces défis comme les principaux freins à l'industrialisation des solutions fondées sur les LLM.

\vspace{0.3cm}

C'est dans ce contexte que s'inscrit notre projet de fin d'études, réalisé au sein de \textbf{TIM Tunisie}, entreprise spécialisée dans l'édition de solutions logicielles pour le secteur industriel. TIM Tunisie développe et commercialise deux systèmes d'information de référence~:

\begin{itemize}[noitemsep]
    \item \textbf{QALITAS}~: plateforme intégrée de gestion de la qualité industrielle~;
    \item \textbf{GMAO PRO}~: système de gestion de la maintenance assistée par ordinateur.
\end{itemize}

Ces deux produits intègrent progressivement des fonctionnalités alimentées par les LLM --- analyse automatique de documents qualité, assistance conversationnelle pour les auditeurs, génération de rapports de conformité --- ce qui génère un besoin croissant de gouvernance et de monitoring de ces ressources d'intelligence artificielle.

\vspace{0.3cm}

Notre mission consiste à concevoir et à développer \textbf{LLM Dashboard}, une plateforme SaaS multi-tenant de gouvernance et de monitoring des modèles de langage. Positionnée comme un \textbf{proxy intelligent} entre les applications clientes et les fournisseurs de LLM, cette plateforme assure~:

\begin{itemize}[noitemsep]
    \item un \textbf{suivi en temps réel} de la consommation de tokens et des coûts associés~;
    \item un \textbf{moteur de gouvernance} fondé sur des règles configurables de quotas, de budgets et de contrôle d'accès~;
    \item un système de \textbf{détection d'anomalies} reposant sur trois algorithmes d'apprentissage automatique complémentaires~;
    \item un module d'\textbf{explication automatique} des anomalies par intelligence artificielle~;
    \item un moteur de \textbf{prévision budgétaire} intégrant une analyse de risques à trois scénarios~;
    \item un système de \textbf{facturation contractuelle} avec génération automatisée de factures~;
    \item un \textbf{tableau de bord interactif} enrichi de visualisations en temps réel.
\end{itemize}

\vspace{0.3cm}

Le développement de cette plateforme a été conduit selon la méthodologie agile \textbf{Scrum}~\cite{schwaber2020}, structuré en trois sprints correspondant aux trois releases fonctionnelles du projet. Cette approche itérative a permis de livrer des incréments fonctionnels réguliers et de recueillir un retour continu de la part de l'entreprise d'accueil.

\vspace{0.3cm}

Le présent rapport est organisé comme suit~:

\begin{itemize}[noitemsep]
    \item Le \textbf{premier chapitre} présente l'étude du projet~: l'organisme d'accueil, la problématique identifiée, la solution proposée et la méthodologie de développement adoptée.
    \item Le \textbf{deuxième chapitre} dresse un état de l'art des technologies mobilisées~: les modèles de langage, la gouvernance des LLM, les architectures multi-tenant et la détection d'anomalies par apprentissage automatique.
    \item Le \textbf{troisième chapitre} détaille le Sprint~0, consacré à l'identification des besoins fonctionnels et non fonctionnels, à la modélisation du cas d'utilisation global, à la définition du product backlog et à l'environnement de travail.
    \item Le \textbf{quatrième chapitre} couvre le Sprint~1, dédié à la mise en place de l'infrastructure technique et des modules fondamentaux~: authentification, gateway LLM, gouvernance et tableau de bord.
    \item Le \textbf{cinquième chapitre} présente le Sprint~2, consacré aux composants d'intelligence artificielle~: détection d'anomalies, explication IA, prévision budgétaire et optimisation des coûts.
    \item Le \textbf{sixième chapitre} couvre le Sprint~3, portant sur les fonctionnalités entreprise --- facturation, traçage des sessions, gestion des prompts --- ainsi que sur le déploiement, les tests et la validation.
\end{itemize}

Le rapport se conclut par un bilan global, une discussion des limites identifiées et une présentation des perspectives d'évolution de la plateforme.
